% BHU PYQ -- Transcribed & Corrected
% Paper: PHIMJ-44 / PHIMJ44: Introduction to Political Philosophy
% Exam: B.A. (Hons.) Semester IV Examination, 2025-26 (NEP)
% Category: Major Core
% Department: Department of Philosophy, Faculty of Arts, Banaras Hindu University

\documentclass[11pt,a4paper]{article}
\usepackage[a4paper,top=2.5cm,bottom=2.5cm,left=2.8cm,right=2.8cm,headheight=14pt]{geometry}
\usepackage{amsmath,amssymb}
\usepackage[T1]{fontenc}
\usepackage[utf8]{inputenc}
\usepackage{lmodern,microtype,fancyhdr,enumitem,hyperref,lastpage,parskip}

\hypersetup{
    colorlinks=true,
    urlcolor=black,
    linkcolor=black,
    pdftitle={PHIMJ44: Introduction to Political Philosophy -- BHU Sem IV 2025-26 (NEP)}
}

\pagestyle{fancy}
\fancyhf{}
\renewcommand{\headrulewidth}{0.4pt}
\renewcommand{\footrulewidth}{0.4pt}
\fancyhead[L]{\small \textit{Banaras Hindu University}}
\fancyhead[R]{\small \textit{PHIMJ-44: Political Philosophy (Sem IV 2025-26)}}
\fancyfoot[C]{\small Page \thepage\ of \pageref{LastPage}}

\newlist{parts}{enumerate}{2}
\setlist[parts,1]{label=(\roman*),leftmargin=*,itemsep=4pt,topsep=2pt}

\newcommand{\pts}[1]{\hfill \textit{[#1]}}

\begin{document}

\begin{center}
  {\large\bfseries BANARAS HINDU UNIVERSITY}\\[2pt]
  {\normalsize B.A. (Hons.) Semester IV Examination, 2025-26 (NEP)}\\[6pt]
  \rule{0.8\linewidth}{0.5pt}\\[6pt]
  {\large\bfseries DEPARTMENT OF PHILOSOPHY}\\[4pt]
  {\large\bfseries Paper Code: PHIMJ-44 : Introduction to Political Philosophy}\\[4pt]
  {\normalsize Time: 2:30 Hours \hfill Max. Marks: 70}
\end{center}

\rule{\linewidth}{0.4pt}

\smallskip

\noindent \textit{(Write your Roll No. at the top immediately on the receipt of this question paper)}\\[2pt]
\noindent \textit{Note: Attempt all the three Sections-A, B and C according to the instructions given.}\\[1pt]
\noindent \textit{निर्देशानुसार अ, ब तथा स तीनों खण्डों के उत्तर दीजिए।}

\bigskip

\begin{center}
  {\large\bfseries SECTION A / खण्ड-अ}\\[2pt]
  {\normalsize (Long Answer Type Questions / दीर्घ उत्तरीय प्रश्न)}
\end{center}

\noindent \textit{Note: Attempt any three questions out of the following five questions, in 300 words. Each question carries 10 marks.}\\[2pt]
\noindent \textit{निम्नलिखित पाँच प्रश्नों में से किन्हीं तीन प्रश्नों का चयन कर प्रत्येक प्रश्न का उत्तर लगभग 300 शब्दों में दीजिए। प्रत्येक प्रश्न 10 अंकों का है।}

\bigskip

\noindent \textbf{Question 1.} Discuss the nature and scope of political philosophy.\\[1pt]
राजनीति-दर्शन के स्वरूप एवं क्षेत्र की विवेचना कीजिए। \pts{10}

\medskip\rule{\linewidth}{0.2pt}\medskip

\noindent \textbf{Question 2.} Elucidate the concept of nationalism.\\[1pt]
राष्ट्रवाद के संप्रत्यय को स्पष्ट कीजिए। \pts{10}

\medskip\rule{\linewidth}{0.2pt}\medskip

\noindent \textbf{Question 3.} Explain the notion of fascism as a political ideology.\\[1pt]
एक राजनीतिक विचारधारा के रूप में फ़ासीवाद की व्याख्या कीजिए। \pts{10}

\medskip\rule{\linewidth}{0.2pt}\medskip

\noindent \textbf{Question 4.} Discuss the concept of equality as a normative principle.\\[1pt]
एक मानकीय सिद्धांत के रूप में ``समानता'' के संप्रत्यय की विवेचना कीजिए। \pts{10}

\medskip\rule{\linewidth}{0.2pt}\medskip

\noindent \textbf{Question 5.} Throw a light on the concept of the state.\\[1pt]
राज्य-सत्ता की अवधारणा पर प्रकाश डालिए। \pts{10}

\bigskip
\rule{\linewidth}{0.2pt}
\bigskip

\begin{center}
  {\large\bfseries SECTION B / खण्ड-ब}\\[2pt]
  {\normalsize (Short Answer Type Questions / लघु उत्तरीय प्रश्न)}
\end{center}

\noindent \textit{Note: Attempt any four questions out of the following six questions, each in 150 words. Each question carries 5 marks.}\\[2pt]
\noindent \textit{निम्नलिखित छः प्रश्नों में से किन्हीं चार प्रश्नों का चयन कर प्रत्येक प्रश्न का उत्तर लगभग 150 शब्दों में दीजिए। प्रत्येक प्रश्न 5 अंकों का है।}

\bigskip

\noindent \textbf{Question 6.} Explain the phenomenon called ``Politics''.\\[1pt]
``राजनीति'' नामक परिघटना की व्याख्या कीजिए। \pts{5}

\medskip\rule{\linewidth}{0.2pt}\medskip

\noindent \textbf{Question 7.} Discuss Aristotle's conception of the State.\\[1pt]
अरस्तू के राज्य-सत्ता के संप्रत्यय की विवेचना कीजिए। \pts{5}

\medskip\rule{\linewidth}{0.2pt}\medskip

\noindent \textbf{Question 8.} Give a brief account of ``freedom'' as a political ideal.\\[1pt]
एक राजनीतिक-आदर्श के रूप में ``स्वतंत्रता'' की व्याख्या कीजिए। \pts{5}

\medskip\rule{\linewidth}{0.2pt}\medskip

\noindent \textbf{Question 9.} What is authority? Explain.\\[1pt]
प्राधिकार क्या है? व्याख्या कीजिए। \pts{5}

\medskip\rule{\linewidth}{0.2pt}\medskip

\noindent \textbf{Question 10.} Discuss in brief the idea of communism.\\[1pt]
साम्यवाद के विचार की संक्षिप्त विवेचना कीजिए। \pts{5}

\medskip\rule{\linewidth}{0.2pt}\medskip

\noindent \textbf{Question 11.} Critically assess Plato's conception of the ideal State.\\[1pt]
प्लेटो के आदर्श-राज्य के संप्रत्यय का आलोचनात्मक मूल्यांकन कीजिए। \pts{5}

\bigskip
\rule{\linewidth}{0.2pt}
\bigskip

\begin{center}
  {\large\bfseries SECTION C / खण्ड-स}\\[2pt]
  {\normalsize (Very Short Answer Type Questions / अतिलघु उत्तरीय प्रश्न)}
\end{center}

\noindent \textit{Note: Attempt all the ten questions of this Section in a maximum of two sentences. Each question carries 2 marks.}\\[2pt]
\noindent \textit{इस खण्ड के सभी दस प्रश्नों का उत्तर अधिकतम दो वाक्यों में दीजिए। प्रत्येक प्रश्न 2 अंकों का है।}

\bigskip

\noindent \textbf{Question 12.}
\begin{parts}
  \item Who defined the State as an association with the monopoly of legitimate violence?\\[1pt]
        किस विचारक ने राज्यसत्ता की परिभाषा ऐसे संघ के रूप में दी जिसके पास हिंसा का वैध एकाधिकार हो? \pts{2}

  \item Who penned ``The Republic''?\\[1pt]
        ``द रिपब्लिक'' नामक ग्रंथ के रचयिता कौन हैं? \pts{2}

  \item What is Sovereignty?\\[1pt]
        संप्रभुता क्या है? \pts{2}

  \item Who considered ``Ethics'' as a part of Politics?\\[1pt]
        किस दार्शनिक ने ``नीतिशास्त्र'' को राजनीति-शास्त्र का भाग माना? \pts{2}

  \item What do you mean by Sarvodaya?\\[1pt]
        सर्वोदय से आपका क्या आशय है? \pts{2}

  \item Who regarded Politics as the supreme human activity?\\[1pt]
        किस विचारक ने राजनीति को सर्वोपरि मानवीय गतिविधि निरूपित किया? \pts{2}

  \item What is liberal democracy?\\[1pt]
        उदारवादी लोकतंत्र क्या है? \pts{2}

  \item Is political philosophy a descriptive science?\\[1pt]
        क्या राजनीति दर्शन एक वर्णनात्मक विज्ञान है? \pts{2}

  \item Who propounded the idea of justice in the soul?\\[1pt]
        आत्मा में निहित न्याय का विचार किस दार्शनिक ने प्रतिपादित किया? \pts{2}

  \item Name a representative philosopher of utopian socialism.\\[1pt]
        स्वप्नदर्शी समाजवाद के किसी प्रतिनिधि दार्शनिक का नाम लिखिए। \pts{2}
\end{parts}

\vfill
\rule{\linewidth}{0.4pt}

\begin{center}\small
\href{https://bhu-exam.vercel.app/}{Click here to visit our website}\\[4pt]
\href{https://me-aryan.vercel.app/}{Click here to visit our contributor}
\end{center}

\end{document}
